\chapter{Results}

\section{Pattern separation}


Some initial results from the protocol described in Section~\ref{sec:separacao_padroes} have already been obtained and will be
discussed in this section.

Figure~\ref{fig:avg_activity} presents the mean activation of GC populations for the control model (without neurogenesis) and
for models with neurogenesis at different iGC-to-EC connectivity levels. In the control model, where all GCs are mature, the
mean population activation, defined as the percentage of cells that spiked at least once during the \SI{1000}{\milli\second}
simulation interval, was 12.46\%, a level of sparse activity, though not as low as that typically observed in the DG. In
future simulations, the impact of increasing inhibitory synapses in the circuit will be analyzed.

The introduction of iGCs into the circuit significantly alters this dynamics. As iGC connectivity increases, their activation
rate rises sharply, consistent with their greater intrinsic excitability. At 100\% connectivity, virtually the entire iGC
population becomes active in response to an input pattern. This increase in iGC activity raises the overall GC population
activation.

In contrast, a progressive decrease in mature Granule Cell (mGC) activity is observed as iGC connectivity increases. This
result suggests the existence of an inhibitory competition mechanism, whereby iGCs, once strongly activated, indirectly suppress
mGC activity through shared interneurons.

However, the suppression observed in mGCs is not as pronounced as that reported in other models~\cite{kimEffect2024}, which may
again indicate that the inhibitory strength in the current model is insufficient. One hypothesis for this discrepancy is the
absence of a direct inhibitory pathway from iGCs to mGCs, a mechanism described by~\citeonline{lunaAdultborn2019} that was not
implemented. The inclusion of this direct connection is a promising avenue for future investigation, as it could confer greater
biological plausibility to the model and will also be explored in the project.

\begin{figure}[H]
    \centering
    \caption{Mean population activation (\%) per model. The vertical axis contains a break and two different scales for better
    visualization. Error is represented by the shaded area.}
    \includegraphics[width=0.7\textwidth]{figuras/plots/avg_activity}
    \label{fig:avg_activity}
\end{figure}


Figure~\ref{fig:pattern_separation} deepens the analysis by detailing the degree of pattern separation ($\mathcal{S}_D$) as a
function of input pattern similarity. A value of $\mathcal{S}_D > 1$ means that the output representation in the DG is more
distinct than the input representation in the EC. In general, all models, including the control, achieve higher pattern
separation for highly similar inputs. This occurs because, mathematically, small differences in output encoding are amplified
when input overlap is very large. The control model's performance with low-similarity inputs was suboptimal, failing to
separate patterns effectively, which can again be attributed to insufficient network sparsity, as discussed earlier. The
introduction of neurogenesis exacerbates this limitation: as iGC connectivity increases, the network's pattern separation
capacity progressively decreases. The model with 100\% iGC connectivity exhibits the worst performance under all similarity
conditions. The underlying mechanism of this deterioration is the non-selective activation of iGCs, which, due to their high
excitability, tend to form the same subset of active neurons in response to distinct input patterns, reducing the dissimilarity
of output representations and, consequently, the degree of pattern separation.

\begin{figure}
    \centering
    \caption{Degree of pattern separation ($\mathcal{S}_D$) per model and input similarity level. The dashed line represents
    a separation degree of 1; values above this indicate that pattern separation occurred, while values below indicate that it
    did not. Error is represented by the shaded area.}
    \includegraphics[width=0.7\textwidth]{figuras/plots/pattern_separation}
    \label{fig:pattern_separation}
\end{figure}

To dissect the individual contributions of each neuronal population, Figure~\ref{fig:avg_pattern_separation} presents the mean
$\mathcal{S}_D$, aggregated across all similarity levels. This analysis confirms the previously observed trend, showing that
the overall GC pattern separation capacity decreases with increasing iGC connectivity. Analysis of the subpopulations reveals a
clear functional dichotomy: iGCs, in isolation, never achieve $\mathcal{S}_D > 1$, indicating that, rather than separating,
they integrate input patterns, making them more similar. In contrast, mGCs demonstrate a robust pattern separation capacity
across all models.

mGC performance improves slightly with increasing iGC connectivity and activity, likely a byproduct of increased global
inhibition in the network, which, though insufficient, may be rendering mGC activity marginally more sparse. A simple linear
regression model was used to evaluate the relationship between iGC connectivity level and mGC degree of pattern separation. The
model found a significant, though weak, relationship between connectivity level and degree of pattern separation
($F(1, 8) = 7.83, p = 0.023, R^2 = 0.49, \mathcal{S}^{mGC}_D = 1.76 + 0.13 \times \text{Connectivity}$).

These findings corroborate the theory proposed by~\citeonline{kimEffect2024}, which suggests a role for mGCs as pattern
separators and for iGCs as integrators. It remains to be investigated, in future experiments, what the impact of this mixed
coding and the strengthening of separation in mGCs has on CA3 memory functions.

\begin{figure}
    \centering
    \caption{Mean degree of pattern separation ($\mathcal{S}_D$) per model and population. The black dashed line represents
    the $\mathcal{S}_D$ of the control, while the gray one represents a separation degree of 1. Error is represented by the
    shaded area; in this case, the error is large given that $\mathcal{S}_D$ varies substantially across different similarity
    levels in Figure~\ref{fig:pattern_separation}.}
    \includegraphics[width=0.7\textwidth]{figuras/plots/avg_pattern_separation}
    \label{fig:avg_pattern_separation}
\end{figure}


\section{Expected results}

It is expected that the presence of iGCs will impair pattern completion in CA3, due to increased activation of neural
assemblies unrelated to the input pattern, an effect consistent with experimental results suggesting non-selective assembly
activation by iGCs due to their high excitability~\cite{koSystems2025}. Additionally, in the temporal maturation simulation,
iGCs are expected to encode patterns over time, integrating information presented while they were young, corroborating their
role as temporal integrators, as proposed by~\cite{aimoneComputational2009}. It is therefore theorized that mGCs would play a
predominant role in the separation of specific, distinct patterns, while iGCs, over the course of their maturation, would
specialize in the separation of patterns close in time, specifically, those they encoded during their immature phase.
